% Options for packages loaded elsewhere
\PassOptionsToPackage{unicode}{hyperref}
\PassOptionsToPackage{hyphens}{url}
\PassOptionsToPackage{dvipsnames,svgnames,x11names}{xcolor}
%
\documentclass[
]{ltjsarticle}

\usepackage{amsmath,amssymb}
\usepackage{iftex}
\ifPDFTeX
  \usepackage[T1]{fontenc}
  \usepackage[utf8]{inputenc}
  \usepackage{textcomp} % provide euro and other symbols
\else % if luatex or xetex
  \usepackage{unicode-math}
  \defaultfontfeatures{Scale=MatchLowercase}
  \defaultfontfeatures[\rmfamily]{Ligatures=TeX,Scale=1}
\fi
\usepackage{lmodern}
\ifPDFTeX\else  
    % xetex/luatex font selection
    \setmainfont[]{UDEV Gothic NF}
\fi
% Use upquote if available, for straight quotes in verbatim environments
\IfFileExists{upquote.sty}{\usepackage{upquote}}{}
\IfFileExists{microtype.sty}{% use microtype if available
  \usepackage[]{microtype}
  \UseMicrotypeSet[protrusion]{basicmath} % disable protrusion for tt fonts
}{}
\makeatletter
\@ifundefined{KOMAClassName}{% if non-KOMA class
  \IfFileExists{parskip.sty}{%
    \usepackage{parskip}
  }{% else
    \setlength{\parindent}{0pt}
    \setlength{\parskip}{6pt plus 2pt minus 1pt}}
}{% if KOMA class
  \KOMAoptions{parskip=half}}
\makeatother
\usepackage{xcolor}
\setlength{\emergencystretch}{3em} % prevent overfull lines
\setcounter{secnumdepth}{5}
% Make \paragraph and \subparagraph free-standing
\makeatletter
\ifx\paragraph\undefined\else
  \let\oldparagraph\paragraph
  \renewcommand{\paragraph}{
    \@ifstar
      \xxxParagraphStar
      \xxxParagraphNoStar
  }
  \newcommand{\xxxParagraphStar}[1]{\oldparagraph*{#1}\mbox{}}
  \newcommand{\xxxParagraphNoStar}[1]{\oldparagraph{#1}\mbox{}}
\fi
\ifx\subparagraph\undefined\else
  \let\oldsubparagraph\subparagraph
  \renewcommand{\subparagraph}{
    \@ifstar
      \xxxSubParagraphStar
      \xxxSubParagraphNoStar
  }
  \newcommand{\xxxSubParagraphStar}[1]{\oldsubparagraph*{#1}\mbox{}}
  \newcommand{\xxxSubParagraphNoStar}[1]{\oldsubparagraph{#1}\mbox{}}
\fi
\makeatother


\providecommand{\tightlist}{%
  \setlength{\itemsep}{0pt}\setlength{\parskip}{0pt}}\usepackage{longtable,booktabs,array}
\usepackage{calc} % for calculating minipage widths
% Correct order of tables after \paragraph or \subparagraph
\usepackage{etoolbox}
\makeatletter
\patchcmd\longtable{\par}{\if@noskipsec\mbox{}\fi\par}{}{}
\makeatother
% Allow footnotes in longtable head/foot
\IfFileExists{footnotehyper.sty}{\usepackage{footnotehyper}}{\usepackage{footnote}}
\makesavenoteenv{longtable}
\usepackage{graphicx}
\makeatletter
\newsavebox\pandoc@box
\newcommand*\pandocbounded[1]{% scales image to fit in text height/width
  \sbox\pandoc@box{#1}%
  \Gscale@div\@tempa{\textheight}{\dimexpr\ht\pandoc@box+\dp\pandoc@box\relax}%
  \Gscale@div\@tempb{\linewidth}{\wd\pandoc@box}%
  \ifdim\@tempb\p@<\@tempa\p@\let\@tempa\@tempb\fi% select the smaller of both
  \ifdim\@tempa\p@<\p@\scalebox{\@tempa}{\usebox\pandoc@box}%
  \else\usebox{\pandoc@box}%
  \fi%
}
% Set default figure placement to htbp
\def\fps@figure{htbp}
\makeatother
% definitions for citeproc citations
\NewDocumentCommand\citeproctext{}{}
\NewDocumentCommand\citeproc{mm}{%
  \begingroup\def\citeproctext{#2}\cite{#1}\endgroup}
\makeatletter
 % allow citations to break across lines
 \let\@cite@ofmt\@firstofone
 % avoid brackets around text for \cite:
 \def\@biblabel#1{}
 \def\@cite#1#2{{#1\if@tempswa , #2\fi}}
\makeatother
\newlength{\cslhangindent}
\setlength{\cslhangindent}{1.5em}
\newlength{\csllabelwidth}
\setlength{\csllabelwidth}{3em}
\newenvironment{CSLReferences}[2] % #1 hanging-indent, #2 entry-spacing
 {\begin{list}{}{%
  \setlength{\itemindent}{0pt}
  \setlength{\leftmargin}{0pt}
  \setlength{\parsep}{0pt}
  % turn on hanging indent if param 1 is 1
  \ifodd #1
   \setlength{\leftmargin}{\cslhangindent}
   \setlength{\itemindent}{-1\cslhangindent}
  \fi
  % set entry spacing
  \setlength{\itemsep}{#2\baselineskip}}}
 {\end{list}}
\usepackage{calc}
\newcommand{\CSLBlock}[1]{\hfill\break\parbox[t]{\linewidth}{\strut\ignorespaces#1\strut}}
\newcommand{\CSLLeftMargin}[1]{\parbox[t]{\csllabelwidth}{\strut#1\strut}}
\newcommand{\CSLRightInline}[1]{\parbox[t]{\linewidth - \csllabelwidth}{\strut#1\strut}}
\newcommand{\CSLIndent}[1]{\hspace{\cslhangindent}#1}

\makeatletter
\@ifpackageloaded{tcolorbox}{}{\usepackage[skins,breakable]{tcolorbox}}
\@ifpackageloaded{fontawesome5}{}{\usepackage{fontawesome5}}
\definecolor{quarto-callout-color}{HTML}{909090}
\definecolor{quarto-callout-note-color}{HTML}{0758E5}
\definecolor{quarto-callout-important-color}{HTML}{CC1914}
\definecolor{quarto-callout-warning-color}{HTML}{EB9113}
\definecolor{quarto-callout-tip-color}{HTML}{00A047}
\definecolor{quarto-callout-caution-color}{HTML}{FC5300}
\definecolor{quarto-callout-color-frame}{HTML}{acacac}
\definecolor{quarto-callout-note-color-frame}{HTML}{4582ec}
\definecolor{quarto-callout-important-color-frame}{HTML}{d9534f}
\definecolor{quarto-callout-warning-color-frame}{HTML}{f0ad4e}
\definecolor{quarto-callout-tip-color-frame}{HTML}{02b875}
\definecolor{quarto-callout-caution-color-frame}{HTML}{fd7e14}
\makeatother
\makeatletter
\@ifpackageloaded{caption}{}{\usepackage{caption}}
\AtBeginDocument{%
\ifdefined\contentsname
  \renewcommand*\contentsname{Table of contents}
\else
  \newcommand\contentsname{Table of contents}
\fi
\ifdefined\listfigurename
  \renewcommand*\listfigurename{List of Figures}
\else
  \newcommand\listfigurename{List of Figures}
\fi
\ifdefined\listtablename
  \renewcommand*\listtablename{List of Tables}
\else
  \newcommand\listtablename{List of Tables}
\fi
\ifdefined\figurename
  \renewcommand*\figurename{図}
\else
  \newcommand\figurename{図}
\fi
\ifdefined\tablename
  \renewcommand*\tablename{Table}
\else
  \newcommand\tablename{Table}
\fi
}
\@ifpackageloaded{float}{}{\usepackage{float}}
\floatstyle{ruled}
\@ifundefined{c@chapter}{\newfloat{codelisting}{h}{lop}}{\newfloat{codelisting}{h}{lop}[chapter]}
\floatname{codelisting}{Listing}
\newcommand*\listoflistings{\listof{codelisting}{List of Listings}}
\makeatother
\makeatletter
\makeatother
\makeatletter
\@ifpackageloaded{caption}{}{\usepackage{caption}}
\@ifpackageloaded{subcaption}{}{\usepackage{subcaption}}
\makeatother

\usepackage{bookmark}

\IfFileExists{xurl.sty}{\usepackage{xurl}}{} % add URL line breaks if available
\urlstyle{same} % disable monospaced font for URLs
\hypersetup{
  pdftitle={ココシルデータ解析０},
  pdfauthor={司馬博文},
  colorlinks=true,
  linkcolor={Cyan3},
  filecolor={Maroon},
  citecolor={Blue},
  urlcolor={Blue},
  pdfcreator={LaTeX via pandoc}}


\title{ココシルデータ解析０}
\usepackage{etoolbox}
\makeatletter
\providecommand{\subtitle}[1]{% add subtitle to \maketitle
  \apptocmd{\@title}{\par {\large #1 \par}}{}{}
}
\makeatother
\subtitle{概観}
\author{司馬博文}
\date{2/26/2025}

\begin{document}
\maketitle


\section{はじめに（1/31/2025
ミーティング内容のまとめ）}\label{ux306fux3058ux3081ux306b1312025-ux30dfux30fcux30c6ux30a3ux30f3ux30b0ux5185ux5bb9ux306eux307eux3068ux3081}

ビリルビンが活性酵素を除去しようとする中で，その残滓であるバイオピリンが活性酸素上昇／酸化ストレスのマーカーになることが知られている．

特に重度の精神疾患患者の判別には AUC 0.779 の精度を持つ (Wake et al.,
2022)．

しかし，職業ストレス (Mourad and Gaballah, 2023)
のほかに，心不全，肝硬変，妊娠，スポーツなど種々の理由でもバイオピリンの尿中排出濃度が変化する．BMI
（特にトリグリセリド）との相関も疑われている．

そこで本研究では，バイオピリンと他変数の依存構造を詳細に検討する．

\begin{tcolorbox}[enhanced jigsaw, leftrule=.75mm, toprule=.15mm, bottomrule=.15mm, breakable, left=2mm, rightrule=.15mm, arc=.35mm, opacityback=0, colback=white, colframe=quarto-callout-tip-color-frame]

\vspace{-3mm}\textbf{リサーチクエスチョン}\vspace{3mm}

\begin{enumerate}
\def\labelenumi{\arabic{enumi}.}
\tightlist
\item
  {究極的には休職・離職を見たい}．昇進との関係などを項目反応理論で補足したい．
\item
  バイピリンと職業性ストレスチェック（57項目）との相関．
\item
  バイオピリンと法定検診との相関：補正や注意事項を構成したい．また一般に他サービスとの連関が得られれば，囲い込みが強くなる．
\end{enumerate}

\end{tcolorbox}

\phantomsection\label{slide-listing}

\subsection{職業性ストレスチェックとの関係}\label{ux8077ux696dux6027ux30b9ux30c8ux30ecux30b9ux30c1ux30a7ux30c3ux30afux3068ux306eux95a2ux4fc2}

第２回 (2/27/2025) ミーティング内容まとめ

\begin{itemize}
\tightlist
\item
  職業性ストレスチェックは産業界で一大サービスであり，それとの関連から説明したい．
\item
  「職業性ストレスチェックのどのような意味での要約になっているか？」を知りたい．
\item
  {弘前の大規模データで検証したい}．
\end{itemize}

\subsection{時系列変化}\label{ux6642ux7cfbux5217ux5909ux5316}

\begin{itemize}
\tightlist
\item
  随時尿にしたからなのか，とんでもない変化率が３ヶ月で出た（７から１０）．

  \begin{itemize}
  \tightlist
  \item
    10月に組織変更があり，職位変更あり群となし群，顧客変更あり群で違いが結構見れる．
  \end{itemize}
\item
  ストレススコアの変化率を調べてもよく判らなかった．
\end{itemize}

\subsection{先行研究}\label{ux5148ux884cux7814ux7a76}

臨床実習のある３回生の方が，４回生よりもバイオピリンレベルが高かった
(Tada et al., 2020)．

(高知恵 et al., 2023)
でも質問票と尿中バイオピリンの双方を用いている．2022年の10月から12月末までのデータである．

しかし「調査当日の受け持ち褥婦数」を調べており，これでは２群間で有意差を認めなかったという結果であった．これには，「翌日から上昇する」という報告を引用して，もう少し時間を空けるべきであった可能性が指摘されている．

\subsection{示唆}\label{ux793aux5506}

\begin{tcolorbox}[enhanced jigsaw, leftrule=.75mm, toprule=.15mm, bottomrule=.15mm, breakable, left=2mm, rightrule=.15mm, arc=.35mm, opacityback=0, colback=white, colframe=quarto-callout-important-color-frame]

\begin{itemize}
\tightlist
\item
  人事に昇進の影響が違うと思われる．やはり \textbf{項目反応モデル}
  だろうか？
\item
  そもそもココシルデータ解析は，「どのようなデータを取れば，有用な結論を導けるか？」というベイズ最適化的な側面があるかもしれない．

  \begin{itemize}
  \tightlist
  \item
    Big Five を取って，関連性が見れたらすごいことになる．
  \item
    NEC
    が出している視線解析のベンチャーがある．尿検査のデータ以外も取れるかもしれない．
  \end{itemize}
\end{itemize}

\end{tcolorbox}

\phantomsection\label{refs}
\begin{CSLReferences}{1}{1}
\bibitem[\citeproctext]{ref-Mourad-Gaballah2023}
Mourad, B. H., and Gaballah, I. F. (2023).
\href{https://journals.lww.com/joem/fulltext/2023/01000/studying_the_association_between_occupational.10.aspx}{Studying
the association between occupational stress and urinary levels of
oxidative stress biomarkers (8-OHdG and biopyrrins) in brickfield
workers}. \emph{Journal of Occupational and Environmental Medicine},
\emph{65}(1).

\bibitem[\citeproctext]{ref-Tada+2020}
Tada, S., Shiota, A., Hayashi, H., and Nakamura, T. (2020).
\href{https://doi.org/10.1016/j.heliyon.2019.e03138}{Reference urinary
biopyrrin level and physiological variation in healthy young adults:
Relation of stress by learning}. \emph{Heliyon}, \emph{6}(1). doi:
10.1016/j.heliyon.2019.e03138.

\bibitem[\citeproctext]{ref-Wake+2022}
Wake, R., Araki, T., Fukushima, M., Matsuda, H., Inagaki, T., Hayashida,
M., \ldots{} Oh-Nishi, A. (2022).
\href{https://doi.org/10.1111/eip.13154}{Urinary biopyrrins and free
immunoglobin light chains are biomarker candidates for screening at-risk
mental state in adolescents}. \emph{Early Intervention in Psychiatry},
\emph{16}(3), 272--280.

\bibitem[\citeproctext]{ref-ux9ad8ux77e5ux6075+2023}
高知恵, 千葉貴子, 中根祥子, 谷口朱子, 北條渉, 中嶋有加里, \ldots{}
渡邊香織. (2023).
\href{https://mol.medicalonline.jp/archive/search?jo=fm0kokus&ye=2023&vo=7&issue=1}{在留外国人褥婦への保健指導を実践する助産師の思いと心理的ストレス
- 尿中バイオピリンと心理尺度による横断的研究 -}.
\emph{国際臨床医学会雑誌}, \emph{7}(1), 22--28.

\end{CSLReferences}




\end{document}
